% Responsáveis: P4 (análise de falhas) + todos (conclusões acionáveis) + P2/P3 (custo)
\section{Discussão}\label{sec:discussao}

\subsection{Análise de falhas}
Como o BERTimbau afinado erra apenas 19 notícias do teste, analisámos todos esses erros e
completámos a amostra com 21 erros do melhor modelo clássico (caracteres com SVM afinado), num total
de 40 (20 de cada tipo). Nove notícias foram mal classificadas pelos dois modelos.
A Tabela~\ref{tab:errors} mostra as categorias. Os falsos negativos (12) são todos notícias
\emph{falsas com estilo jornalístico}: datas, cargos entre parênteses, instituições e fontes citadas,
como a notícia falsa sobre o depoimento de Lula à Polícia Federal (E10). Os falsos positivos
dividem-se entre notícias verdadeiras \emph{com tom de opinião ou citações} (8), como uma reportagem
sobre terraplanistas que começa só com as palavras deles (E02), e notícias verdadeiras sobre
\emph{temas típicos das falsas} (7), como polémicas de celebridades ou Lula e a Lava Jato (E03, E08).

Os dois modelos erram de forma diferente. O modelo clássico falha sobretudo por pistas
superficiais: o tema (6 de 21) e o estilo jornalístico imitado (8 de 21). O BERTimbau, que capta
melhor o tom, falha sobretudo em textos verdadeiros de opinião ou análise (5 dos 10 erros só
seus), que associa às notícias falsas. Entre os 9 erros comuns, 3 parecem \emph{erros de rótulo
no corpus}: notícias reais copiadas por sítios de desinformação, algumas ainda com a assinatura da
agência ou do jornal de origem (``agências/R7'', ``O Tempo''). No total, 5 dos 40 erros têm
rótulo duvidoso, o que sugere que parte do ``erro'' restante do BERTimbau é ruído do próprio corpus,
rotulado pelo sítio de origem e não pela veracidade de cada notícia. Por fim, 4 erros vêm de textos
quase vazios (9 a 20 palavras no corpus), sem informação suficiente para decidir.
\todo{P4: rever este comentário (categorias de E10--E40 sugeridas pelo assistente)}

\begin{table}[t]\centering\scriptsize
\caption{Categorias de erro do melhor modelo no teste (40 erros analisados). FP: verdadeira classificada como falsa; FN: falsa classificada como verdadeira.}\label{tab:errors}
\begin{tabular}{lrrrl}
\toprule
Categoria & FP & FN & Total & Exemplos \\
\midrule
Falsa com estilo jornalístico & 0 & 12 & 12 & E09, E10, E15 \\
Verdadeira com tom de opinião ou citações & 8 & 0 & 8 & E02, E06, E13 \\
Tema típico da outra classe & 7 & 0 & 7 & E03, E04, E08 \\
Possível erro de rótulo no corpus & 0 & 5 & 5 & E07, E12, E14 \\
Pouco texto informativo no corte & 1 & 3 & 4 & E01, E05, E26 \\
Outro & 3 & 0 & 3 & E17, E25, E33 \\
Marcas da fonte & 1 & 0 & 1 & E18 \\
\bottomrule
\end{tabular}
\end{table}


\subsection{Resposta à pergunta da organização}
Os três tipos de análise apontam na mesma direção: no Fake.Br, o que distingue as notícias falsas
é \emph{como} são escritas (frases curtas, registo informal, erros, ausência de datas
concretos), e não \emph{de que} falam. Daqui retiramos três conclusões acionáveis:
\begin{enumerate}
\item \textbf{Não confiar em métricas que o comprimento explica.} Qualquer avaliação de um detetor
deve incluir a \emph{baseline} de comprimento; no Fake.Br ela atinge 0,944.
\item \textbf{Usar a ausência de datas como sinal de triagem.} As verdadeiras têm cerca de cinco
vezes mais datas concretas, confirmado com anotação humana. Notícias sem nenhuma data devem subir na
fila de verificação. É um sinal simples, explicável ao verificador e menos dependente do tema; um
extrator de entidades genérico não basta para os outros tipos, porque confunde instituições com locais.
\item \textbf{Usar o BERTimbau afinado como detetor principal e os n-gramas de caracteres como
alternativa barata.} O primeiro é o mais exato mesmo sem atalhos (0,987); o segundo corre em CPU,
generaliza melhor do que as palavras para temas novos e explica-se com os atributos mais pesados.
Em ambos os casos, reavaliar periodicamente, porque o desempenho em temas não vistos cai até 5,5 pontos.
\end{enumerate}
As três conclusões mantêm-se no teste (Tabela~\ref{tab:test}).

\subsection{Custo computacional}
A Tabela~\ref{tab:cost} resume o custo dos modelos finais, medido com o mesmo protocolo
(\texttt{src/cost.py}) num computador com Windows~11, processador AMD de 16 núcleos lógicos e 15\,GB de
RAM. Os modelos clássicos treinam em 10 a 43\,s, ocupam no máximo 3,5\,MB e classificam uma notícia
em 1 a 11\,ms, com menos de 200\,MB de memória. O extrator de entidades (\texttt{md}) processa cada
notícia em cerca de 36\,ms; o modelo \texttt{lg} tem mais 577\,MB do que o \texttt{md} para um
ganho de 1,6 pontos de F1 (Secção~\ref{sec:metodo}). O BERTimbau precisa de GPU: no Colab gratuito
(T4), o \emph{fine-tuning} com 5 folds demorou 12 minutos com 3,6\,GB de memória, e o treino final
mais a previsão das 1\,444 notícias de teste 167\,s; o modelo ocupa cerca de 436\,MB. Os cinco
pontos de F1 que o BERTimbau ganha aos modelos clássicos custam, portanto, uma GPU e um modelo cerca
de cem vezes maior. O pipeline não usa modelos de linguagem nem serviços externos pagos (0 tokens,
0 euros); o único recurso externo foi a GPU gratuita do Colab.

\begin{table}[t]\centering\scriptsize
\caption{Custo computacional. Treino no dev completo (5\,756 notícias); inferência medida em 1\,000 notícias. Memória: pico de RAM do processo (GPU no caso do BERTimbau). CPU: Windows 11, 16 núcleos lógicos, 14,9 GB de RAM; GPU: Google Colab gratuito. Nenhum serviço externo pago nem modelo de linguagem foi usado no pipeline (0 tokens).}\label{tab:cost}
\begin{tabular}{llrrrr}
\toprule
Modelo & Hardware & Treino (s) & ms/notícia & Memória (MB) & Tam. (MB) \\
\midrule
Só comprimento & CPU & 0,0 & 0,00 & 0 & 0,0 \\
palavras + SVM & CPU & 10,1 & 1,45 & 110 & 3,5 \\
caracteres + SVM & CPU & 39,5 & 10,83 & 129 & 1,5 \\
caracteres + LogReg & CPU & 42,7 & 10,86 & 175 & 3,0 \\
spaCy (md) & CPU & -- & 36,44 & 284 & 54,1 \\
BERTimbau afinado: CV 5 folds & GPU T4 & 713,2 & -- & 3587 & 436,0 \\
BERTimbau afinado: final + teste & GPU T4 & 167,4 & -- & 3580 & 436,0 \\
\bottomrule
\end{tabular}
\end{table}

